\section{Depth and Practical Value}
\label{sec:depth}

\paragraph{Depth profiles.} For each of nine input corruptions, a model's sensitivity
fingerprint is the relative change of its residual stream at each layer; two models are compared by the rank
correlation of their fingerprints on a shared relative-depth axis; activations are never moved between models. For
TimesFM and Chronos-2 the agreement replicates on the v2 private split for 6 of 9 corruptions (overall $\rho=0.524$
[0.504, 0.544]; Figure~\ref{fig:heldout}a) and for 4 of 9 on v3, with 8 of 9 keeping their sign. The two models react at
matching depths to detrending (private $\rho=0.368$) and dropout (0.510), and at opposite depths to noise ($-0.743$),
spikes ($-0.785$) and warps ($-0.651$). Frequency shift (0.381 against 0.216 on dev) and smoothing ($-0.581$ against
$-0.418$) fail the registered rule because the agreement is stronger on private data than on dev. Only for level shift
is the agreement unstable; it changes sign between dev and private on both corpora (\fid{CA-01}).

\paragraph{Window-level restoration.} Patching clean activations window by window into a
corrupted TimesFM pass shows that the forecast is restored mainly in the window just before the horizon, increasingly so with
depth (from 0.427 of the damage undone at layer 0 to 0.612 at layer 8 and 0.631 at layer 16 for noise), while restoration
in the first window falls from $-0.191$ to 0.000 (320 series; \fid{DE-01}). This recency effect is
not visible when whole layers are patched.

\paragraph{Comparability of depth.} The captured blocks account for 13.9\%
(Chronos-T5-Base, encoder only) to 99.8\% (Timer) of a forecast's FLOPs (\fid{DE-05}), and on a whole-stack depth axis
Chronos-T5's last encoder block sits at depth 0.478 (\fid{DE-04}). Which layers appear most important also depends
strongly on the selection criterion (\fid{DE-08}). We therefore compare depth profiles across models rather than
the locations of single layers.

\paragraph{Failure prediction.} For each model we predict a series' log MASE from
free, output-only features (quantile width, forecast flatness, catch22 features of the context;
\citealp{lubba2019catch22}), add dictionary-family activations, and measure the gain, cross-fitted with a series
bootstrap. The free baseline is strong (Spearman 0.76--0.85), and for five models the gain's interval includes zero. The
gain is positive for the two least accurate models, Timer (+0.030 [0.019, 0.043]) and Time-MoE (+0.012 [0.005,
0.019]), and both confirm on the private split (Timer 0.0446 [0.0292, 0.0605], $p_{\text{Holm}}=0.002$; Time-MoE 0.011
[0.0032, 0.0194], $p_{\text{Holm}}=0.012$) and again on v3 (\fid{MN-28}). Internals do not help in choosing a model. Routing
each series to the model with the lowest predicted MASE gives 1.964 [1.837, 2.093] with internals against 1.948 [1.814,
2.070] without, both worse than always using Chronos-2 (1.856). For the more accurate models, the
model's output carries as much information about its errors as its internal features do.
